For \(r\in\{0,1\}\), let \(\Phi^{(r)}(y)\) denote the candidate target expression obtained by inserting \(q_C^{(r)}\) in the \(C\)-factor and retaining the common outer conditional \(p^S\) and the common factors \(q_j\), \(D_j\ne C\).  By \(\mathrm{def:eq23\text{-}functional}\), all state spaces and therefore all displayed sums are finite, and
\[
 \Phi^{(r)}(y)
 =\sum_{w,d_{D\setminus Y_N}}
 p^S(y_S,w\mid x_S)
 q_C^{(r)}
 \prod_{D_j\ne C}q_j .
\]
Subtracting the formula for \(r=1\) from that for \(r=0\), and using distributivity over the finite sums, gives
\[
 \Phi^{(0)}(y)-\Phi^{(1)}(y)
 =\sum_{w,d_{D\setminus Y_N}}
 p^S(y_S,w\mid x_S)
 \bigl(q_C^{(0)}-q_C^{(1)}\bigr)
 \prod_{D_j\ne C}q_j
 =\Delta(y).
\]
Two probability measures on the finite state space \(\prod_{v\in Y}\mathcal X_v\) are unequal if and only if they differ on at least one singleton.  Consequently the two constructed targets differ if and only if \(\Delta(y)\ne0\) for at least one outcome value \(y\).

It remains to verify the asserted obstruction.  Consider two strictly positive kernels on binary coordinates \((A_0,R_0)\), defined by
\[
 \begin{aligned}
 q_C^{(0)}(0,0)&=q_C^{(0)}(1,1)=\frac38,&
 q_C^{(0)}(1,0)&=q_C^{(0)}(0,1)=\frac18,\\
 q_C^{(1)}(a,r)&=\frac14&&\text{for every }(a,r)\in\{0,1\}^2.
 \end{aligned}
\]
Each table sums to one, every entry is positive, and the tables are unequal.  Nevertheless, for each \(r\in\{0,1\}\),
\[
 \sum_{a\in\{0,1\}}
 \bigl(q_C^{(0)}(a,r)-q_C^{(1)}(a,r)\bigr)=0.
\]
Take \(A_0\) to be among the coordinates summed over in \(D\setminus Y_N\), and take the common positive factor
\(p^S(y_S,w\mid x_S)\prod_{D_j\ne C}q_j\) to be independent of \(A_0\).  Summing first over \(A_0\) makes every summand of the remaining contraction zero, so \(\Delta(y)=0\) for every \(y\).  Thus strict positivity and \(q_C^{(0)}\ne q_C^{(1)}\) do not imply nonvanishing of the finite contraction.